Our results show that the RK4 method with 32000 steps over a $50 \mu s$ can simulate an ideal Penning trap to a final, relative error of $2.31 \cdot 10^{-11}$ in comparison to the analytical solution derived in \cref{sec:penningtrap}.
Furthermore, we found empirical evidence for the global error of the RK4 method to be of order $\mathcal{O}(dt^4)$, which is in line with the theoretical order. In comparison, the less accurate FE method showed a final, relative error of 0.06. 
The introduction of a second particle into the system led to significant coupling between the two particles, as they transferred energy through the Coulomb force. We can clearly see this in Figures \ref{fig:2-particle} and \ref{fig:phase-space-plots}, where the simulated paths and phase spaces of the two particles are plotted.\\
The subsequent analysis of the perturbed case demonstrated the phenomenon of resonance in Penning traps. We conducted an inquiry into the influence of perturbation amplitude and angular frequency on particle loss, identifying the ranges where the effect is most pronounced (around resonant frequencies). We've also shown that an increase in amplitude clearly correlates with the widening of said ranges. The following higher-resolution sweeps showed that including inter-particle forces in the simulation yields noticeably different, noisier results.
Ultimately, these resonance effects expose the inherent instability of the trap under specific time-dependent perturbations.
An applied frequency near the system's resonance leads to the rapid escape of the confined particles, resulting in a complete breakdown of the trap's underlying function. \\
A possible usage of a more specialized numerical method, such as a leapfrog/Verlet algorithm in combination with a Boris push, could lead to more physically sound results\cite{LectureNotes}. This could be especially useful for a larger simulation period, where a more significant energy drift in the RK4 method could occur. \\
The analysis of perturbed Penning traps, although successful, was quite computationally demanding. This necessitated limiting the simulation to 100 particles and choosing a larger time step, which could negatively impact accuracy. The introduction of multithreading or GPU-parallelisation could greatly accelerate these computations, as simulations for separate amplitudes and frequencies constitute independent problems and are therefore trivially parallelisable.
